\documentclass[11pt,a4paper]{article}
\usepackage[utf8]{inputenc}
\usepackage{amsmath,amssymb,amsfonts}
\usepackage{geometry}
\geometry{margin=1in}

\title{MOSFETQ DOS / oeneyeOS: Release Candidate v0.4 Specification \\ \large High-Assurance Memory Boundaries, Tier-Certified Addressing, and Structural Wave Mechanics}
\author{Kai Olaf Ketelhut}
\date{September 2026}

\begin{document}

\maketitle

\begin{abstract}
\noindent \textbf{Scope:} This document establishes Release Candidate v0.4 ($\text{rc v0.4}$) for the MOSFETQ DOS runtime architecture and SDK v23 ecosystem. \\
\textbf{Assurance Target:} It defines the formal specification and EAL7-aligned memory protection boundaries, integrates Tier-certified 48-bit virtual address safety envelopes, and formalizes structural wave mechanics within $D_5$ manifolds to govern deterministic state transitions.
\end{abstract}

\section{System Versioning Matrix}
\begin{table}[h]
\centering
\begin{tabular}{|l|l|l|}
\hline
\textbf{Layer} & \textbf{Identifier} & \textbf{Status} \\ \hline
Runtime Core & v20.0 & Active Execution Engine \\ \hline
Kernel Core & v0.2 & Pinned, Immutable Floppy Architecture \\ \hline
Developer Kit & SDK v23 & Active Build Environment \\ \hline
Specification & \textsc{rc v0.4} & High-Assurance Architecture Draft \\ \hline
\end{tabular}
\end{table}

\section{EAL7 Core Memory Protection Boundaries}
To satisfy formal verification requirements for high-assurance environments:
\begin{enumerate}
    \item \textbf{Isolation Invariant:} The pinned v0.2 kernel core address space remains strictly disjoint from user-space execution contexts ($\text{W:}$ workspace and mounted container drives).
    \item \textbf{State Transition Atomicity:} All memory mutations are modeled as total functions ensuring zero side effects on unmapped sectors, supported by formal refinement proofs (comparable to Coq/Isabelle-verified seL4 models).
    \item \textbf{Privilege Ring Enforcement:} Ring 0 (kernel primitives) and Ring 3 (SDK v23 userland/NT emulation layers) are separated by hard hardware-enforced boundary checks.
\end{enumerate}

\section{Tier-Certified Hardware Addressing Envelope}
For fault-tolerant, high-availability deployments, logical address allocation adheres to strict safety thresholds:
\begin{itemize}
    \item \textbf{Virtual Address Ceiling:} Constrained to a deliberate 48-bit hardware safety envelope ($2^{48} - 1$ addresses, representing a 256 TB byte capacity ceiling), preventing translation drift across redundant nodes.
    \item \textbf{Unified Memory Partitioning:} Upper memory segments are reserved for mirrored, failover-safe unified address spaces to eliminate split-brain states during live migration.
\end{itemize}

\section{Structural Wave Mechanics ($D_5$ Manifold Framework)}
To replace probabilistic wave functions with a deterministic structural formalism for state transitions (and framed as an open physics proposal requiring ongoing manifold and operator verification):
\begin{equation}
    \mathcal{E}_{\text{dissipation}} = \oint_{\partial D_5} \left( \mathbf{T}_{\text{structural}} \cdot \mathbf{I}_{\text{inclusion}} \right) d\Sigma
\end{equation}
\begin{itemize}
    \item \textbf{Lattice Relaxation:} State transitions and energy transfers are modeled as mechanical consequences of lattice adjustments within $D_5$ manifolds.
    \item \textbf{Offset Tensors:} Complex-valued wave functions are superseded by structural offset tensors and inclusion operators to resolve discrete binary topological constraints.
    \item \textbf{Quaternion Partitioning:} Modular quaternion partitions act as Doppler tensor anchors, ensuring stable, singularity-free energy dissipation across memory state transitions.
\end{itemize}

\end{document}
